\section{Reproducibility and Implementation Details}
\label{app:reproducibility}

The reproducible unit of the framework is a dataset-construction round. A round starts from a dataset goal and a set of object concepts, then produces prompts, white-background object images, foreground masks, reconstructed meshes, scene-aware rendered views, review records, train-ready pairs, a downstream checkpoint, and a comparison report. These artifacts form a linked state trace that lets later inspection connect a benchmark regression to the samples and construction stages that produced it.

The case-study dataset uses a source-target pair schema. A source image is the canonical front view, and each target image corresponds to a requested horizontal rotation. Each row records the source image, target image, instruction text, target rotation, object identifier, object name, and prompt version. The instruction format uses explicit view names, for example asking the model to rotate a named object from the front view to a right-side or back view.

Split construction is object-disjoint. The semantic-object prompt baseline uses 35 train objects, 7 validation objects, and 8 test objects, yielding 245 train pairs, 49 validation pairs, and 56 test pairs. The first feedback-guided scaling configuration, internally labeled R1, adds 20 new train-only objects and 60 weak-angle pairs while leaving validation and test objects fixed. This design prevents new synthetic objects from leaking into held-out splits and makes the comparison a targeted train-data intervention.

The training probe uses \qwenedit{} with LoRA fine-tuning under a fixed comparison protocol \cite{hu2021lora,wu2025qwenimage,qwen2025qwenimageedit2511}. The tracked runs use rank 32, learning rate \(1\times10^{-4}\), 30 epochs, and the epoch 29 checkpoint for evaluation. The probe is evaluated on \spatialedit{} \cite{xiao2026spatialedit}, which contains 488 rotation pairs across 61 objects and eight angle slots. The metric set combines traditional image metrics, representation metrics, distribution-level quality, and VLM-based view/consistency scores, matching the interpretation used in \cref{sec:results}.

\begin{table}[t]
\centering
\small
\caption{Render-quality diagnosis for the first feedback-guided scaling round, used to route the round to \inspect{}.}
\label{tab:r1-quality-diagnosis}
% Auto-generated by arxiv_report/paper/scripts/generate_latex_assets.py
\begin{tabular}{ll}
\toprule
Diagnostic item & Value \\
\midrule
New feedback-scaling objects & 20 \\
Objects below hybrid-score gate 0.6 & 20 \\
Median latest hybrid score & 0.447 \\
Min / max latest hybrid score & 0.365 / 0.582 \\
Verdict & inspect \\
Next action & render-quality gating before merge \\
\bottomrule
\end{tabular}

\end{table}

\begin{table}[t]
\centering
\small
\caption{Augmentation composition for the first feedback-guided scaling round.}
\label{tab:r1-augmentation}
% Auto-generated by arxiv_report/paper/scripts/generate_latex_assets.py
\begin{tabular}{lrl}
\toprule
Group & Count & Unit \\
\midrule
daily\_item & 10 & objects \\
sports\_item & 3 & objects \\
street\_object & 4 & objects \\
vehicle & 3 & objects \\
target yaw 90$^\circ$ & 20 & pairs \\
target yaw 180$^\circ$ & 20 & pairs \\
target yaw 270$^\circ$ & 20 & pairs \\
\bottomrule
\end{tabular}

\end{table}

Checkpoint and artifact availability remains partial in this technical-report package. The tracked comparison uses the \semanticprobe{} baseline checkpoint, internally labeled exp5, and the \feedbackprobe{} R1 checkpoint, but the large checkpoint files and full inference image dumps are not bundled because of size. A complete release should attach qualitative grids, pipeline diagrams, loss curves, failure-case figures, and checkpoint MD5 or SHA-256 hashes.

\paragraph{Artifact index.}
Code, intermediate artifacts, and evaluation logs will be released at an anonymous repository upon publication. Key hyperparameters are LoRA rank \(=32\), learning rate \(=10^{-4}\), and training epoch \(=29\), held uniform across the compared rounds. The dataset split is approximately 305 train pairs, 49 validation pairs, and 56 test pairs after the feedback-guided augmentation. Checkpoint hashes will be provided in the released repository.
